\documentclass[border=4pt]{standalone}
\usepackage{amsmath,amssymb,mathtools,xcolor,varwidth}
\begin{document}
\begin{varwidth}{28em}
\large\bfseries
Here is the whole proof in miniature. Had no force acted at $B$, inertia (\textbf{Law I}) would carry the      body straight on to $c$, with $Bc = AB$; then triangles $SAB$ and $SBc$ share the base      $SB$ and have equal heights, so \emph{$SAB = SBc$}. But at $B$ a single centripetal impulse toward      $S$ deflects the body, and by the parallelogram of motions (\textbf{Law II, Cor.~1}) it lands at $C$ where      $Cc \parallel SB$. Because $Cc$ is parallel to $SB$, triangles $SBc$ and $SBC$ stand on the same      base $SB$ between the same parallels — so \emph{$SBc = SBC$}, and therefore \emph{$SBC = SAB$}. This is      Newton's genius: an impulse pointed straight at $S$ bends the body's \emph{direction} but leaves the swept      triangular area untouched.
\end{varwidth}
\end{document}
